\subsection{Introduction}

\subsubsection{Security}

\indent 

\textbf{Security} is the practice of protecting information systems including hardware, software, networks, and data, from unauthorized access, misuse, disruption, or destruction.

3 dimensions of security in IT:

\begin{itemize}
    \item Physical Security: Protecting the tangible infrastructure that supports IT. Goal: prevent theft, damage, or physical tampering with assets.
    \item Digital/Information Security: Protecting data and digital systems from cyber threats. Goal: maintain confidentiality, integrity, and availability of data.
    \item Operational Security: Protecting systems through processes, people, and policies. Goal: reduce risks caused by human error, insider threats, or weak processes.
\end{itemize}

\subsubsection{Security Management}

\indent 

Security Management is the process of identifying, implementing, monitoring, and continuously improving security measures to protect an organization’s assets, both physical and digital.

It ensures that risks are managed systematically, and security aligns with business objectives, compliance requirements, and operational needs.

Main goals of Security management:

\begin{itemize}
    \item Protect Confidentiality, Integrity, and Availability (CIA).
    \item Ensure business continuity.
    \item Maintain compliance and legal obligations.
    \item Build trust with users, customers, and stakeholders.
\end{itemize}

\bigskip

\noindent \textbf{Key elements of Security Management}

\begin{itemize}
    \item Risk Assessment and Analysis: Identify threats, vulnerabilities, and potential impacts.
    \item Policy and Governance: Define security policies, standards, and compliance frameworks (ISO 27001, NIST, GDPR, etc.).
    \item Implementation of Controls
    \item Monitoring and Detection: Continuous monitoring of systems, networks, and behaviors to detect anomalies.
    \item Incident Response and Recovery: Plans and procedures for containment, eradication, recovery, and lessons learned.
    \item Continuous Improvement: Regular audits, penetration testing, red/blue team exercises, feedback loops.
\end{itemize}

\bigskip 

\noindent\textbf{Why does Security management matter in IT?}

\begin{itemize}
    \item Protection of Critical Assets: IT systems store intellectual property, customer data, financial records, and operational information. Without proper security management, these assets are at risk of theft, manipulation, or destruction.
    \item Mitigating Cyber Threats: Organizations face constant attacks: malware, phishing, ransomware, insider threats. Security management ensures proactive defenses and rapid responses.
    \item Ensuring Business Continuity: Outages or breaches can cripple operations. Strong security management minimizes downtime, enabling resilience during attacks or disasters.
    \item Legal and Regulatory Compliance: Laws like GDPR, HIPAA, and standards like ISO 27001 mandate security practices. Failure to comply leads to fines, lawsuits, and reputational harm.
    \item Safeguarding Reputation and Trust: Customers, partners, and stakeholders expect organizations to handle data responsibly. Breaches erode trust and can cause lasting brand damage.
    \item Managing Human Risk: People are often the weakest link (phishing, weak passwords, negligence). Security management addresses this with training, policies, and monitoring.
    \item Cost Avoidance: The cost of prevention is far lower than the financial and operational impact of a breach.
\end{itemize}

\subsection{Understanding Vulnerabilities}

\subsubsection{Factors contributing to Vulnerabilities}

\indent 

\begin{itemize}
    \item Human Factors: Weak passwords, password reuse. Falling for phishing/social engineering. Lack of awareness or training. Insider threats (malicious or negligent employees).
    \item Technological Factors: Unpatched software and outdated systems. Misconfigurations (firewalls, cloud storage, databases). Weak encryption or lack of encryption. Zero-day vulnerabilities in hardware/software.
    \item Organisational Factors: Lack of clear policies and governance. Poor incident response planning. Insufficient monitoring and auditing. Over-reliance on third-party vendors without proper vetting.
    \item Physical/Environmental Factors: Insecure server rooms or data centers. Theft or tampering with devices. Natural disasters (floods, fires, earthquakes). Power failures or inadequate redundancy.
    \item Process and Operational Factors: Inadequate access control (too many privileges). Poor patch/change management processes.
\end{itemize}

\subsubsection{The biggest security threat - Users}

\indent

Despite all technical protections, humans are often the weakest link in security. Even the strongest firewalls, encryption, and monitoring can be bypassed if users make mistakes or act carelessly.

\noindent\textbf{Why?}

\begin{itemize}
    \item Human error: Sending sensitive information to the wrong recipient or Misconfiguring systems or accidentally deleting data.
    \item Poor password practices: Weak passwords, password reuse across accounts. Sharing credentials with others.
    \item Phishing and Social engineering: Falling for emails, calls, or messages designed to steal credentials. Clicking malicious links or downloading infected attachments.
    \item Insider threats: Disgruntled employees intentionally misusing access.
\end{itemize}

\subsection{Cyber-Crimes}

\indent 

Cyber-crimes are criminal activities that involve the internet, computers, 
networks, or digital devices. They are actions that violate laws by targeting information systems, data, or digital communications for malicious purposes. Cyber-crimes can target individuals, organizations, or governments. They often exploit technical vulnerabilities or human weaknesses. Prevention requires a combination of technical controls, policies, and user awareness.

\subsubsection{Identity Theft and Fraud}

\indent 

Identity theft occurs when someone steals another person’s personal information (like name, Social Security/Tax ID, credit card, or login credentials) to commit fraudulent activities, such as financial theft, unauthorized transactions, or impersonation.

Preventative measures:

\begin{itemize}
    \item Use strong, unique passwords and enable multi-factor authentication (MFA).
    \item Monitor accounts for suspicious activity regularly.
    \item Be cautious with sharing personal information online or over email/phone.
    \item Shred sensitive documents before disposal.
    \item Install antivirus and anti-malware software on all devices.
\end{itemize}

 \subsubsection{Harassment and Cyber-bullying}

 \indent 

\noindent\textbf{Harassment}: Repeated, unwanted behavior intended to intimidate, threaten, or disturb an individual.

\noindent\textbf{Cyberbullying}: A form of harassment that occurs \textbf{online or through digital devices}, targeting individuals with malicious intent.

Both involve misuse of digital platforms to \textbf{cause emotional, psychological, or reputational harm}.

Preventative Measures:

\begin{itemize}
    \item Implement \textbf{acceptable use policies} and online conduct rules.
    \item Monitor platforms and enforce \textbf{reporting mechanisms}.
    \item Train users on recognizing and responding to harassment.
    \item Legal and regulatory frameworks (e.g., cyber harassment laws in Australia)
\end{itemize}

\subsubsection{Hacking and Unauthorised access}

\indent

Hacking / Unauthorized access = gaining access to systems, networks, data, or accounts without permission, or using legitimate access in an unauthorized way. It covers anything from a simple password guess to sophisticated intrusion campaigns.

Preventative measures:

\begin{itemize}
    \item Strong password policies + mandatory multi-factor authentication (MFA).
    \item Least-privilege access and role-based access control (RBAC).
    \item Regular security awareness training (phishing simulations).
    \item Vendor/third-party risk assessments.
    \item Web Application Firewall (WAF) and input validation.
    \item Encryption for data at rest and in transit.
\end{itemize}

\subsubsection{Data Breaches}

\indent

A data breach occurs when sensitive, protected, or confidential information is accessed, disclosed, or stolen without authorization. This can involve personal data (PII), financial records, trade secrets, or intellectual property.

Preventative Measures:

\begin{itemize}
    \item Encrypt data at rest and in transit
    \item Implement strong access controls (least privilege, RBAC)
    \item Regular security awareness training
    \item Patch and vulnerability management
    \item Monitor and audit systems continuously
    \item Backup critical data and test recovery plans
\end{itemize}

\subsubsection{Malware}

\indent 

\textbf{Malware} (malicious software) is any software intentionally designed to \textbf{damage, disrupt, steal, or gain unauthorized access} to computer systems, networks, or data.

Preventative measures:

\begin{itemize}
    \item Install and update antivirus/anti-malware software
    \item Keep systems and software patched and up-to-date
    \item Educate users on phishing and safe browsing practices
    \item Use firewalls and intrusion detection systems
    \item Backup critical data regularly and securely
    \item Restrict unnecessary administrative privileges
\end{itemize}

\subsubsection{Phishing and Social Engineering}

\indent

\noindent\textbf{Social Engineering}: Manipulating or deceiving people to reveal confidential information or perform actions that compromise security.

\noindent\textbf{Phishing}: A type of social engineering that uses \textbf{emails, messages, or websites} to trick users into giving sensitive data (e.g., passwords, credit cards).

The key idea is that the attackers exploit \textbf{human psychology}, not just technical vulnerabilities.

Preventative measures:

\begin{itemize}
    \item User awareness and training: teach staff to recognize suspicious emails, links, and messages.
    \item Multi-factor authentication (MFA): prevents stolen credentials from being enough to access accounts.
    \item Email and web filters: block malicious content before it reaches users.
\end{itemize}

\subsubsection{DDoS Attacks}

\indent 

A \textbf{Distributed Denial of Service (DDoS)} attack aims to \textbf{overwhelm a target (server, service, network)} with traffic or requests from many distributed sources so that legitimate users cannot access it. The goal is \textbf{disruption} rather than data theft (though DDoS is sometimes used as a smokescreen for other attacks).

Preventative measures:

\begin{itemize}
    \item \textbf{Traffic scrubbing / scrubbing centres} (third-party DDoS mitigation services).
    \item \textbf{Rate limiting} and upstream filtering with ISP cooperation.
    \item \textbf{Autoscaling with careful limits}: helps absorb spikes but can be costly and must be coupled with filtering.
    \item \textbf{Load balancers} and redundancy across regions.
    \item Centralized logging, SIEM, and real-time traffic monitoring
\end{itemize}

\subsubsection{Advanced Persistent Threats (APTs)}

\indent 

An \textbf{Advanced Persistent Threat (APT)} is a \textbf{prolonged and targeted cyberattack} in which an intruder gains unauthorized access to a network and remains undetected for an extended period. APTs usually target \textbf{high-value assets} like intellectual property, financial data, or government secrets.

\begin{itemize}
    \item Preventative measures:
    \item Continuous network monitoring, anomaly detection, endpoint detection.
    \item Close known security gaps that APTs might exploit.
    \item A detailed plan to detect, contain, and remediate intrusions.
    \item Deploying Least privilege, network segmentation, and multi-factor authentication
\end{itemize}

\subsection{Security Threats and Organisational Challenges}

\indent 

\subsubsection{Categories of Security Threats}

\indent 

Security threats are potential events or actions that can compromise the confidentiality, integrity, or availability of information or systems. They can be intentional, accidental, or natural.

There are mainly 5 categories of security threats:

\begin{itemize}
    \item Human based: Threats caused by \textbf{people}, either deliberately or unintentionally.
    \item Technical: Threats exploiting \textbf{hardware, software, or network vulnerabilities}.
    \item Physical: Threats targeting the \textbf{tangible infrastructure} that supports IT.
    \item Operational: Threats arising from \textbf{weak processes, policies, or operational failures}.
    \item External: Threats from \textbf{outside the organization’s direct control}.
\end{itemize}

\subsubsection{Factors making Organisations vulnerable}

\indent 

Organizations are vulnerable to attacks and breaches due to a combination of technical, human, operational, and environmental factors. Understanding these helps in effective security management.

\begin{itemize}
    \item Human factors:
        \begin{itemize}
            \item - Lack of security awareness/training → employees fall for phishing or social engineering.
            \item Weak password practices → reuse, simplicity, or sharing credentials.
            \item Negligence or mistakes → misconfigured systems, accidental deletion.
        \end{itemize}
    \item Technological factors:
        \begin{itemize}
            \item Outdated software or unpatched systems → exploits are available for attackers.
            \item Use of legacy systems → incompatible with modern security standards.
        \end{itemize}
    \item Operational Factors:
        \begin{itemize}
            \item Weak access control policies → excessive privileges, lack of role-based access.
            \item Inadequate incident response and recovery plans.
        \end{itemize}
    \item Physical Factors:
        \begin{itemize}
            \item Unprotected facilities → theft, tampering, or natural disasters.
            \item Poor disaster recovery planning → inability to recover after an incident.
        \end{itemize}
    \item External Factors:
        \begin{itemize}
            \item Third-party/vendor vulnerabilities → compromised suppliers can affect your systems.
            \item Regulatory or legal changes → unprepared organizations may be non-compliant.
            \item Market pressures → cutting corners on security due to cost or time constraints.
        \end{itemize}
\end{itemize}

\subsubsection{Factors that make Security Harder}

\indent 

Security in IT isn’t just about technology.. it’s complex, dynamic, and multidimensional. 

Several factors make protecting systems and data increasingly difficult.

\begin{itemize}
    \item Rapid technological changes
    \item Increasing sophistication of attacks
    \item Human factors
    \item Complexity of IT environments
    \item Regulatory and Compliance Pressures
    \item Resource constraints
    \item Evolving threat landscape
\end{itemize}

\subsection{Information Security}

\indent 

Information security (often shortened to InfoSec) is the practice of \textbf{protecting information}, in any form (digital, physical, or spoken), from \textbf{unauthorized access, use, disclosure, disruption, modification, or destruction}.

At its core, it ensures that data remains CIA:

\begin{itemize}
    \item \textbf{Confidential}: Only authorized people or systems can access it.
    \item \textbf{Integrity-protected}: Information is accurate, reliable, and not altered improperly.
    \item \textbf{Available}: Information and systems are accessible when needed.
\end{itemize}

\subsubsection{CIA Triad}

\indent 

The CIA Triad is the foundational model for information security. It focuses on three core principles that every organization should protect:

\begin{itemize}
    \item Confidentiality: Ensuring that information is accessible only to authorized individuals or systems. Prevents unauthorized access or disclosure. Techniques include encryption, access controls and permissions, and multi-factor authentication.
    \item Integrity: Ensuring that information is accurate, complete, and unaltered except by authorized actions. Prevents unauthorized modification or tampering. Techniques include checksums, hashes, audit logs and version control, data validation, etc.
    \item Availability: Ensuring that information and systems are accessible and usable when needed by authorized users. Prevents downtime, service disruptions, and data loss. Techniques include backups, load balancing, etc.
\end{itemize}

\subsubsection{How threats impact CIA triad?}

\indent 

Security threats can compromise Confidentiality, Integrity, or Availability, the three pillars of the CIA Triad. Understanding these impacts helps organisations design better controls.

\begin{itemize}
    \item Confidentiality threats: Unauthorized access or disclosure of sensitive data.
    \item Integrity threats: Unauthorized modification, corruption, or deletion of data.
    \item Availability threats: Systems or data become unavailable to authorized users.
\end{itemize}

\subsubsection{Domains of Information Security}

\begin{itemize}
    \item Cybersecurity: Protecting digital systems, networks, and applications from unauthorized access, disruption, or attack.
    \item Physical Security: Protecting the tangible assets that support information systems (people, hardware, facilities).
    \item Operational Security: Protecting processes, workflows, and organizational practices to ensure information remains safe.
    \item Data Security: Protecting the information itself, whether it’s at rest, in transit, or in use.
\end{itemize}

\subsubsection{Why Information Security matters for IT professionals?}

\begin{itemize}
    \item Protecting Sensitive Data: IT professionals are custodians of critical organisational data (financial records, intellectual property, personal data, client information). A breach can cause financial loss, reputational damage, and legal consequences.
    \item Enabling Trust and Reputation: Organizations thrive on trust from customers, partners, and employees. IT professionals are on the front lines of building systems that are secure by design. A single security lapse (e.g., data leak, ransomware attack) can undermine years of trust.
    \item Legal and Regulatory Compliance: Governments enforce data protection and privacy laws such as GDPR (EU), HIPAA (US), Privacy Act (Aus), etc. IT professionals must design and maintain systems that comply with these frameworks, or the organization risks lawsuits and heavy penalties.
    \item Supporting Business Continuity: Cyberattacks, insider threats, and natural disasters can disrupt services. IT professionals ensure availability of systems through disaster recovery, backups, and incident response. Without security planning, downtime = loss of productivity, revenue, and customer loyalty.
    \item Defending against increasing Cyberthreats: Threat actors are more sophisticated, organized, and well-funded than ever. IT pros must stay ahead of threats like phishing, social engineering, malware, ransomware, DDoS attacks, etc. Security skills are essential in every IT role, not just for “security specialists.”
    \item Professional Responsibility and Ethics: IT professionals often have privileged access to sensitive data. Ethical handling of data is part of their duty to employers, clients, and society. Misuse or negligence can cause personal liability and career damage.
\end{itemize}

\subsubsection{Individual Rights of Data}

\indent 

In the digital age, people have specific rights over their \textbf{personal data} (any information that identifies them, for example name, email, biometrics, financial info). 

These rights are protected by \textbf{laws (GDPR, CCPA, Privacy Act 1988 in 
Australia)} and reflect broader \textbf{ethical principles}.

\begin{itemize}
    \item Right to be informed: Individuals must know \textbf{how their data is collected, stored, used, and shared}. Organizations are required to provide \textbf{clear privacy notices} (not hidden in fine print).
    \item Right of access: Individuals can request access to the \textbf{personal data an organization holds about them}. IT systems must support easy and timely \textbf{data access requests}.
    \item Right to rectification (correction): If data is \textbf{inaccurate or incomplete}, individuals have the right to correct it.
    \item Right to erasure (right to be forgotten): Individuals can request their personal data be deleted if:
        \begin{itemize}
            \item It’s no longer needed for the purpose it was collected.
            \item They withdraw consent.
            \item It was collected unlawfully.
        \end{itemize}
        Exception: Data needed for legal, contractual, or public interest purposes.
    \item Right to restrict processing: Individuals can limit how their data is processed. Example: Pausing marketing emails while keeping an account active.
    \item Right to data portability: Individuals can request their data in a \textbf{structured, machine-readable format} (e.g., CSV, JSON). This enables them to move data between providers (e.g., switching banks, telecoms)
    \item Rights related to automated decision making and profiling: Individuals can object to: \textbf{Direct marketing} (ads, promotions), Processing based on \textbf{legitimate interests} or \textbf{public interest tasks}. Organizations must stop processing unless they prove \textbf{compelling legitimate grounds}.
\end{itemize}

\subsection{Questions}

\subsubsection{End of Lecture Questions}

\indent 

\begin{itemize}
    \item Define security in the context of IT and explain its three main dimensions.
    \item What is security management and why is it important in IT organizations?
    \item List at least five factors that contribute to organizational vulnerabilities.
    \item Describe the CIA Triad and give a practical example for each element.
    \item Differentiate between phishing and social engineering.
    \item Explain the difference between a DDoS attack and an Advanced Persistent Threat (APT).
\end{itemize}

\subsubsection{Case Study}

\indent

A mid-sized healthcare organization recently suffered a cyberattack. An employee clicked a phishing email link, which installed ransomware on several hospital servers. This caused temporary unavailability of patient records, and some sensitive patient data was also exposed. The attackers demanded payment to decrypt files.

\bigskip

\noindent\textbf{Case Study based Questions}

\begin{itemize}
    \item Identify which CIA Triad elements were impacted and explain how.
    \item Discuss the human, technological, and operational factors that contributed to this breach.
    \item Propose at least five preventive measures that could have mitigated this attack.
    \item What incident response steps should the organization take immediately after detecting the attack?
    \item Which legal and ethical considerations must the organization address post-attack?
\end{itemize}